\chapter*{ВСТУП}
\addcontentsline{toc}{chapter}{ВСТУП}
\thispagestyle{fancy}

\textbf{Актуальність теми.}
В умовах зростаючих загроз кібербезпеці та широкого впровадження дистанційної роботи
захищений доступ до корпоративних ресурсів набуває критичного значення. Технологія
VPN є основним засобом забезпечення конфіденційності та цілісності мережевого трафіку.
Проте традиційні VPN-рішення часто не забезпечують належного рівня автентифікації
користувачів, обмежуючись лише перевіркою криптографічних ключів. Відсутність
двофакторної автентифікації (2FA) залишає такі системи вразливими до атак
на основі компрометації облікових даних. Сучасний протокол WireGuard, незважаючи
на свої переваги у продуктивності та безпеці, за своєю природою не підтримує
механізмів автентифікації користувачів. Таким чином, розробка системи, яка поєднує
WireGuard з надійним механізмом 2FA, є актуальною науково-технічною задачею.

\textbf{Мета і завдання дослідження.}
Метою дипломної роботи є розробка VPN-застосунку з двофакторною автентифікацією
на основі протоколу WireGuard. Для досягнення мети поставлено такі завдання:
\begin{enumerate}
  \item проаналізувати існуючі технології VPN та методи двофакторної автентифікації;
  \item дослідити архітектуру та криптографічні основи протоколу WireGuard;
  \item провести порівняльний аналіз існуючих рішень інтеграції автентифікації з WireGuard;
  \item розробити архітектуру системи VPN з двофакторною автентифікацією;
  \item реалізувати програмний комплекс та провести його тестування.
\end{enumerate}

\textbf{Об'єкт та предмет дослідження.}
Об'єктом дослідження є процес автентифікації користувачів у VPN-мережах.
Предметом дослідження є методи та засоби інтеграції двофакторної автентифікації
з протоколом WireGuard.

\textbf{Методи дослідження.}
У роботі використано такі методи дослідження: порівняльний аналіз протоколів
та методів автентифікації; метод системного аналізу при проектуванні архітектури;
методи об'єктно-орієнтованого програмування при реалізації; методи тестування
програмного забезпечення при верифікації результатів.

\textbf{Наукова новизна та практичне значення результатів.}
Наукова новизна роботи полягає в тому, що удосконалено підхід до інтеграції
двофакторної автентифікації з протоколом WireGuard шляхом застосування моделі
керування життєвим циклом пірів, що, на відміну від існуючих рішень, не потребує
модифікації протоколу WireGuard та використання попередньо розділених ключів (PSK).
Практичне значення роботи полягає у розробці функціонального VPN-застосунку,
який реалізує механізм двофакторної автентифікації на основі стандарту TOTP
(RFC~6238) і підтримує управління ключами WireGuard через захищений API.

\textbf{Структура та обсяг роботи.}
Дипломна робота складається зі вступу, трьох розділів, висновків, списку
використаних джерел та додатків. У першому розділі проводиться аналіз
технологій VPN, криптографічних основ протоколу WireGuard, методів двофакторної
автентифікації та існуючих рішень інтеграції автентифікації з WireGuard.
У другому розділі описується проєктування архітектури системи. У третьому розділі
викладається реалізація програмного комплексу та наведено результати тестування системи.

\newpage
